\documentclass[../readings.tex]{subfiles}

\begin{document}

\subsection{Introduction to Norms}
\label{sub:norms-intro}

A \emph{norm} is a function that assigns a non-negative scalar ``size'' to a vector or
matrix. Norms are the fundamental tool for measuring errors, bounding perturbations,
and defining convergence throughout scientific computing. This chapter introduces the
key vector norm definitions and the families most used in practice. We will study norms
more rigorously in later chapters, where we develop matrix norms, operator
inequalities, Cauchy-Schwarz, Hölder's inequality, and the condition number.

\addDef{Vector Norm}{%
A function $\|\cdot\| : \fR^n \to \fR$ is a \emph{vector norm} if for all
$\bx, \by \in \fR^n$ and $c \in \fR$:
\begin{enumerate}
  \item $\|\bx\| \geq 0$ and $\|\bx\| = 0 \Leftrightarrow \bx = \bzero$ \quad (positive definiteness)
  \item $\|c\bx\| = |c|\|\bx\|$ \quad (absolute homogeneity)
  \item $\|\bx + \by\| \leq \|\bx\| + \|\by\|$ \quad (triangle inequality)
\end{enumerate}
}
\label{def:vector-norm-intro}

\addDef{$\ell^p$ Norms}{%
For $\bx \in \fR^n$ and $p \geq 1$, the \emph{$\ell^p$ norm} is
\[
  \|\bx\|_p = \left(\sum_{i=1}^n |x_i|^p\right)^{1/p}.
\]
The three most common cases are:
\begin{itemize}
  \item $\|\bx\|_1 = \displaystyle\sum_{i=1}^n |x_i|$ \quad (\emph{Manhattan norm} --- sum of absolute values)
  \item $\|\bx\|_2 = \sqrt{\displaystyle\sum_{i=1}^n x_i^2} = \sqrt{\bx^T\bx}$ \quad (\emph{Euclidean norm})
  \item $\|\bx\|_\infty = \max_{1 \leq i \leq n} |x_i|$ \quad (\emph{max norm}, the limit as $p \to \infty$)
\end{itemize}
}
\label{def:p-norms-intro}

\addNote{%
\textbf{Physical Intuition.}\enspace
In scientific computing, the choice of norm depends on what you want to measure:
\begin{itemize}
    \item \textbf{$\ell_2$ (Energy):} The standard measure of magnitude; corresponds to physical energy or distance.
    \item \textbf{$\ell_\infty$ (Worst-case):} Measures the single largest entry. Useful for ``tolerance'' checks (e.g., ``every entry must be below $10^{-6}$'').
    \item \textbf{$\ell_1$ (Sparsity):} Encourages or measures sparsity; often used in optimization and compressed sensing.
\end{itemize}
}

\addExmpl{%
For $\bx = (3, -4, 0)^T$:
\[
  \|\bx\|_1 = 3 + 4 + 0 = 7, \qquad
  \|\bx\|_2 = \sqrt{9 + 16 + 0} = 5, \qquad
  \|\bx\|_\infty = 4.
\]
The unit balls $\{\bx : \|\bx\|_p \leq 1\}$ in $\fR^2$ are: a diamond ($p = 1$), a
circle ($p = 2$), and a square ($p = \infty$). As $p$ increases, the unit ball grows
toward the $\ell^\infty$ square.
}

\addNote{%
Throughout these notes, $\|\cdot\|$ without a subscript denotes the Euclidean norm
$\|\cdot\|_2$ unless stated otherwise. Errors and residuals are almost always measured
in $\|\cdot\|_2$. In NumPy: \texttt{np.linalg.norm(x)} computes $\|\bx\|_2$;
\texttt{np.linalg.norm(x, ord=p)} computes $\|\bx\|_p$.
}

\addNote{%
\textbf{Numerical Zero.}\enspace
In floating-point arithmetic, the condition $\|\bx\| = 0$ is rarely met exactly due to rounding. We instead use a \textbf{tolerance} $\varepsilon$: we say $\bx$ is ``numerically zero'' if $\|\bx\| \leq \varepsilon \cdot \|\bx_0\|$, where $\bx_0$ is some initial reference vector.
}

\addThm{All Norms are Equivalent on $\fR^n$}{%
For any two norms $\|\cdot\|_\alpha$ and $\|\cdot\|_\beta$ on $\fR^n$, there exist
constants $c_1, c_2 > 0$ such that
\[
  c_1\|\bx\|_\beta \leq \|\bx\|_\alpha \leq c_2\|\bx\|_\beta \quad \text{for all } \bx \in \fR^n.
\]
In particular: $\|\bx\|_2 \leq \|\bx\|_1 \leq \sqrt{n}\,\|\bx\|_2$ and
$\|\bx\|_\infty \leq \|\bx\|_2 \leq \sqrt{n}\,\|\bx\|_\infty$.
}
\label{thm:norm-equivalence-intro}

\addPrf{%
Since $\fR^n$ is finite-dimensional, the unit sphere $\{\bx : \|\bx\|_\beta = 1\}$ is
compact. Any norm is a continuous function, so it attains its minimum and maximum on a
compact set --- these extremal values are $c_1$ and $1/c_2$. The explicit inequalities
follow from Cauchy-Schwarz: $\|\bx\|_1 = \sum|x_i| \cdot 1 \leq \|\bx\|_2\,\|\mathbf{1}\|_2 = \sqrt{n}\|\bx\|_2$.
}

\addExcr{%
Refer to Defs~\ref{def:vector-norm-intro}--\ref{def:p-norms-intro} and
Thm~\ref{thm:norm-equivalence-intro}.
\begin{enumerate}
  \item For $\bx = (-1, 2, -3, 4)^T$, compute $\|\bx\|_1$, $\|\bx\|_2$, $\|\bx\|_\infty$
  by hand. Verify all three norm-equivalence inequalities hold.
  \item Verify the triangle inequality $\|\bx + \by\|_2 \leq \|\bx\|_2 + \|\by\|_2$
  for $\bx = (1, 2)^T$, $\by = (-3, 1)^T$ by computing both sides.
  \item Using NumPy, confirm your hand calculations with
  \texttt{np.linalg.norm(x, 1)}, \texttt{np.linalg.norm(x)}, and
  \texttt{np.linalg.norm(x, np.inf)}.
  \item Show that $\|\bx\|_\infty = \lim_{p \to \infty}\|\bx\|_p$.
  (\textit{Hint: let $M = \max_i|x_i|$; factor $M$ out of $\|\bx\|_p$ and take the limit.})
\end{enumerate}
}

\end{document}
